%% chapters/01-pendahuluan.tex
%% VERSI GELAR-STYLE — Struktur, gaya narasi, dan panjang mengikuti
%% disertasi Gelar Budiman (audio watermarking, ITB).
%% Seksi: I.1 Latar Belakang | I.2 Masalah Penelitian | I.3 Tujuan
%%        I.4 Kontribusi dan Kebaruan | I.5 Metodologi Penelitian
%% Referensi dipertahankan dari LaTeX asli.

\chapter{Pendahuluan}
\label{bab:pendahuluan}

Bagian pendahuluan ini menjelaskan latar belakang, rumusan
masalah, tujuan penelitian, \textcolor{blue}{manfaat penelitian,} kontribusi dan kebaruan penelitian, \textcolor{blue}{metodologi penelitian, batasan dan asumsi, serta sistematika
penulisan disertasi}.


%% ================================================================
\section{Latar Belakang}
\label{sec:bab1_latar}
%% ================================================================

Pemerintah Indonesia  berkomitmen untuk membangun industri manufaktur
yang kuat dengan meluncurkan \emph{Making Indonesia 4.0} pada April 2018,
dengan tujuan meningkatkan produktivitas sektor manufaktur dan daya saing
ekspor global \citetitb{Kemenperin2018Making4}.
Program ini menetapkan lima sektor industri prioritas yaitu makanan dan minuman,
tekstil, kimia, otomotif, dan elektronika, sebagai ujung tombak penerapan
teknologi industri 4.0, yang mencakup \emph{Internet of Things} (IoT),
analisis \emph{big data}, \emph{advanced }robotika, dan kecerdasan artifisial
(\emph{artificial intelligence}, AI) \citetitb{WEF2018Industry4}.
Salah satu implementasi industri 4.0 yang berdampak langsung pada efisiensi
rantai produksi adalah perawatan cerdas atau \emph{smart maintenance},
yang mengandalkan data sensor dan kecerdasan artifisial untuk memprediksi
kegagalan komponen sebelum hal tersebut terjadi.

Saat ini terdapat dua jenis perawatan yang umum diterapkan pada industri
manufaktur yaitu perawatan korektif (\emph{corrective maintenance}) dan perawatan
preventif (\emph{preventive maintenance}). 
Perawatan korektif dilakukan setelah kerusakan terjadi, sehingga
mengakibatkan produksi berhenti secara tidak terduga dan menimbulkan kerugian
finansial yang signifikan.
Perawatan preventif berupaya mengatasi permasalah kegagalan dengan mengganti bagian komponen atau mesin setelah
jangka waktu tertentu \citetitb{Mobley2002PlantMaintenance}.
Namun pendekatan ini memiliki kelemahan fundamental, di mana komponen yang masih
layak pakai diganti lebih awal sehingga menambah biaya, atau bahkan sebaliknya
komponen yang sudah mulai rusak atau aus justru terlewat dari proses inspeksi.
Teknologi industri 4.0, khususnya kecerdasan artifisial, \emph{machine learning},
dan analisis \emph{big data}, memungkinkan lahirnya paradigma perawatan baru
yang bersifat prediktif, di mana intervensi dilakukan berdasarkan kondisi
aktual komponen yang diukur secara waktu nyata \emph{real-time}, bukan berdasarkan mekanisme terjadwal
yang tetap \citetitb{Lei2018Machinery}.

Di antara banyaknya komponen mesin rotasi, \emph{rolling element bearing}
berada di posisi yang paling kritis sekaligus paling rentan terhadap kegagalan.
\emph{Bearing} digunakan secara luas pada hampir seluruh mesin rotasi
industri, mulai dari motor listrik, pompa, kompresor, hingga lini otomasi
manufaktur \citetitb{Lei2018Machinery}.
Studi empiris menunjukkan bahwa kegagalan \emph{bearing} menyumbang sekitar
40--50\,\% dari total kegagalan mesin rotasi, menjadikan hal tersebut penyebab utama
\emph{downtime} yang tidak terencana di industri manufaktur
\citetitb{Nectoux2012PRONOSTIA}.
Gambar~\ref{fig:bab1_objek_industri} memperlihatkan komponen \emph{rolling
element bearing} beserta tahapan proses \textcolor{blue}{permesinannya}, yang menjadi konteks dalam penelitian ini.
Pada penelitian ini, konteks industri bersumber dari PT~SKF~Indonesia yang
mengoperasikan mesin \emph{grinding outer ring} (OR1 dan OR2) untuk
memproduksi cincin luar \emph{bearing} bagi sektor otomotif
\citetitb{SKF2017TechReport}.
Apabila perawatan prediktif dapat diterapkan secara efektif pada proses
permesinan \emph{bearing}, maka waktu henti atau kegagalan yang tidak terencana dapat
diturunkan, masa pakai alat komponen dapat dioptimalkan, dan masalah kualitas
produk akibat alat komponen yang aus dapat ditekan secara sistematis.

\begin{figure}[ht]
  \centering
  \begin{minipage}[b]{0.42\textwidth}
    \centering
    \includegraphics[width=\linewidth]{ball_bearing_components.png}
    \caption*{(a) Komponen \emph{rolling element bearing}}
  \end{minipage}\hfill
  \begin{minipage}[b]{0.55\textwidth}
    \centering
    \includegraphics[width=\linewidth]{ball_bearing_machining.png}
    \caption*{(b) Proses permesinan \emph{outer ring} dan \emph{inner ring}}
  \end{minipage}
  \caption{Konteks fisik objek penelitian: (a)~komponen utama
           \emph{rolling element bearing}, dan (b)~tahapan \emph{grinding}
           dan \emph{honing} pada cincin luar dan dalam yang menjadi sumber
           sinyal vibrasi mesin produksi.
           Sumber: SKF Group~\citetitb{SKF2017TechReport}.}
  \label{fig:bab1_objek_industri}
\end{figure}

Kemampuan mendeteksi jenis kerusakan \emph{bearing}, apakah terjadi pada
\emph{inner race}, \emph{outer race}, atau bola, maupun memprediksi sisa
umur pakai (\emph{remaining useful life}, RUL) secara akurat, merupakan dua
pilar utama sistem perawatan prediktif berbasis kondisi yang efektif.
Pendekatan deteksi dan prediksi degradasi \emph{bearing} mengalami perkembangan
pesat dalam satu dekade terakhir, dari model fisika berbasis persamaan
keausan Hertzian sampai pada metode berbasis data yang memanfaatkan rekaman
sinyal degradasi berskala besar \citetitb{Lei2018Machinery}.

Pada dekade terakhir, penggunaan arsitektur \emph{deep learning} mendominasi literatur
kasus diagnostik dan prognostik pada \emph{bearing}.
Untuk kasus diagnostik, berbagai arsitektur jaringan saraf konvolusi \emph{convolution neural networks} telah
mencapai akurasi sangat tinggi pada dataset benchmark CWRU
(\emph{Case Western Reserve University}) \citetitb{Zhang2017WDCNN}.
Sebagai contoh, arsitektur \emph{Wide Deep Convolutional Neural Network}
(WDCNN) \citetitb{Zhang2017WDCNN} mampu mengklasifikasikan sepuluh kelas
kerusakan langsung dari sinyal vibrasi mentah tanpa memerlukan ekstraksi
fitur secara manual.
Model berbasis \emph{Long Short-Term Memory} (LSTM) \citetitb{Hochreiter1997LSTM}
dan Transformer \citetitb{Vaswani2017Attention} memperkuat kemampuan  
ke arah pemodelan sekuens temporal untuk estimasi RUL.
Arsitektur terbaru seperti Mamba \citetitb{Gu2023Mamba} dan xLSTM
\citetitb{Beck2024xLSTM} membahas pemodelan sekuens yang efisien secara komputasi
yang lebih baik untuk kebutuhan prognostik di lingkungan produksi nyata.

Namun, tingginya akurasi model \emph{deep learning} tidak mendorong tingginya adopsi teknologi tersebut di industri manufaktur.
Tantangan utamanya adalah kurangnya interpretabilitas (\emph{interpretability})
dari keputusan model: insinyur perawatan dan pemeliharaan tidak dapat menerima prediksi dari
sistem kotak hitam untuk mengambil keputusan bernilai tinggi, seperti
kapan harus menghentikan mesin atau mengganti komponen tanpa dapat memverifikasi
alasan mendasar secara fisik dari prediksi tersebut \citetitb{Mobley2002PlantMaintenance}.
Permasalahan ini kemudian mendorong berkembangnya bidang \emph{Explainable
Artificial Intelligence} (XAI), yang bertujuan membuat model kecerdasan
artifisial dapat diaudit, dijelaskan, diinterpretasi, lalu dipercaya oleh pengguna yang notabenennya tidak memiliki kemampuan teknis.

Penerapan XAI pada diagnostik \emph{bearing} selama ini hampir seluruhnya
beroperasi di tingkat fitur agregat, seperti kurtosis, energi pita
frekuensi, atau nilai RMS, yang merupakan turunan dari sinyal mentah (\emph{raw signal}).
Pendekatan ini tidak mampu menjawab pertanyaan mendasar tentang mekanisme
pembelajaran model, di bagian mana dari sinyal mentah model memusatkan
perhatiannya, dan apakah bagian tersebut sesuai dengan impuls transien
fisik yang dikenal dari konsep kegagalan dan kerusakan \emph{bearing}?
Di kasus prognostik, model prediksi RUL berbasis \emph{deep learning}
menghasilkan estimasi yang akurat, namun representasi laten modelnya tetap tidak
transparan, belum ada bukti statistik bahwa model-model tersebut
menginternalisasi degradasi \emph{bearing} secara fisika, khususnya frekuensi
karakteristik BPFO (\emph{Ball Pass Frequency Outer race}), BPFI
(\emph{Ball Pass Frequency Inner race}), BSF (\emph{Ball Spin Frequency}),
dan FTF (\emph{Fundamental Train Frequency}), di dalam representasi laten
yang dipelajari.

Berdasarkan latar belakang tersebut, disertasi ini mengusulkan pendekatan
perawatan prediktif yang mencakup \textcolor{blue}{dua hal} yang saling mendukung dan dapat
dikombinasikan, yaitu \textcolor{blue}{(i)}~diagnostik kerusakan \emph{bearing} berbasis
\emph{benchmark} multi-model dengan interpretabilitas dari sinyal mentah, dan
\textcolor{blue}{(ii)}~prognostik RUL dengan interpretabilitas mekanistik berbasis
\emph{Sparse Autoencoder} yang dipetakan ke frekuensi karakteristik
degradasi \emph{bearing}.
\textcolor{blue}{Kedua metode tersebut} dirancang agar dapat ditempatkan dalam arsitektur komputasi
bertingkat, dari sensor di tepi jaringan hingga \emph{cloud}, sehingga
sistem yang dihasilkan tidak hanya akurat, tetapi juga dapat diinterpretasikan, dan dapat
dipercaya oleh praktisi industri.
Skema pendekatan penelitian disertasi ini diilustrasikan pada
Gambar~\ref{fig:bab1_skema}.

\begin{figure}[ht]
  \centering
  \includegraphics[width=0.88\textwidth]{research_flow.png}
  \caption{Skema pendekatan penelitian: dari tantangan perawatan prediktif
           industri menuju kerangka diagnostik--prognostik yang akurat,
           terjelaskan, dan dapat diadopsi di lingkungan manufaktur.}
  \label{fig:bab1_skema}
\end{figure}


%% ================================================================
\section{Rumusan Penelitian}
\label{sec:bab1_rumusan}
%% ================================================================

Pada sistem perawatan prediktif \emph{bearing} berbasis \emph{deep learning},
terdapat tiga permasalahan utama yang saling berkaitan.
\textcolor{blue}{Pertama}, model dengan akurasi tinggi cenderung kompleks dan tidak dapat diinterpretasikan,
sementara model yang sederhana dan mudah dijelaskan sering kali kurang
akurat pada kasus kerusakan yang bervariasi.
\textcolor{blue}{Kedua}, model yang mampu \textcolor{blue}{mendeteksi kerusakan dan} memprediksi sisa umur pakai secara akurat
tidak selalu dapat membuktikan bahwa \textcolor{blue}{keputusan dan} prediksinya berlandaskan pada fisika
degradasi yang sebenarnya terjadi, bukan sekadar korelasi statistik \textcolor{blue}{semata}.
\textcolor{blue}{Ketiga, model yang dikembangkan dan divalidasi pada dataset benchmark
laboratorium belum tentu dapat bekerja dengan baik pada data industri nyata,
sehingga kelayakan penerapannya untuk mengurangi waktu henti yang tidak terencana
pada proses produksi \emph{ball bearing} belum terbukti.}
Berdasarkan kajian literatur yang telah dilakukan, dapat dirumuskan
masalah penelitian sebagai berikut:

\begin{enumerate}

  \item \textcolor{blue}{Bagaimana mengembangkan model perawatan prediktif \textcolor{blue}{\emph{ball bearing}}
          berbasis data sensor vibrasi dengan pendekatan \emph{machine learning}
          dan \emph{deep learning} yang mampu melakukan deteksi kerusakan
          (\emph{fault diagnosis}) dan prediksi sisa umur pakai
          (\emph{remaining useful life}) secara akurat?}

  \item \textcolor{blue}{Bagaimana menerapkan metode \emph{explainable artificial intelligence}
        (XAI) pada model perawatan prediktif \textcolor{blue}{\emph{ball bearing}}
        berbasis data sensor vibrasi berbasis
        \emph{deep learning}, sehingga keputusan klasifikasi kerusakan dapat
        menghasilkan peta karakteristik kerusakan diskriminatif serta
        representasi laten model prediksi RUL dapat dibuktikan bersesuaian
        dengan fisika degradasi \emph{bearing}?}

  \item \textcolor{blue}{Bagaimana menerapkan dan memvalidasi model prediksi
        sisa umur pakai (\emph{remaining useful life}) \emph{bearing} yang
        dikembangkan pada dataset benchmark terhadap data sensor vibrasi dari
        proses produksi \emph{ball bearing} di industri manufaktur real?}

\end{enumerate}


% \begin{enumerate}

%   \item Bagaimana mengembangkan model perawatan prediktif \textcolor{blue}{\emph{bearing}}
%         berbasis data sensor vibrasi dengan pendekatan \emph{machine learning}
%         dan \emph{deep learning} yang mampu melakukan deteksi kerusakan
%         (\emph{fault diagnosis}) dan prediksi sisa umur pakai
%         (\emph{remaining useful life}) secara akurat, sehingga
%         mendukung pengurangan waktu henti yang tidak terencana pada sistem
%         produksi \textcolor{blue}{\emph{ball bearing}}?

%   \item \textcolor{blue}{Bagaimana menerapkan metode \emph{explainable artificial intelligence}
%         (XAI) pada model diagnostik dan prognostik \emph{bearing} berbasis
%         \emph{deep learning}, sehingga keputusan klasifikasi kerusakan dapat
%         menghasilkan peta karakteristik kerusakan (\emph{fault signature})
%         yang stabil dan diskriminatif pada level sinyal mentah, serta
%         representasi laten model prediksi RUL dapat dibuktikan bersesuaian
%         dengan fisika degradasi \emph{bearing} (frekuensi karakteristik BPFO,
%         BPFI, BSF, dan FTF), sehingga hasil model dapat diverifikasi dan
%         dipercaya oleh insinyur perawatan?}

%   \item \textcolor{blue}{Bagaimana menerapkan dan memvalidasi model prediksi
%         sisa umur pakai (\emph{remaining useful life}) \emph{bearing} yang
%         dikembangkan pada dataset benchmark terhadap data sensor vibrasi dari
%         proses produksi \emph{ball bearing} di industri manufaktur real,
%         serta merancang arsitektur sistem perawatan prediktif multi-tier yang
%         mendukung pengurangan waktu henti yang tidak terencana dan menjadi
%         fondasi bagi adopsi di lingkungan produksi nyata?}

% \end{enumerate}

Berdasarkan rumusan masalah di atas, dikumpulkan beberapa referensi
yang merupakan \emph{state of the art} (SoTA) dari bidang yang terkait.
Referensi-referensi telah membahas banyak keunggulan, 
namun masih memiliki celah yang dapat dikembangkan, sehingga terdapat
ruang bagi kebaruan dan kontribusi dalam bidang terkait.

\textbf{Diagnostik \emph{bearing} berbasis sinyal mentah.}
\citenameitb{Zhang2017WDCNN} mengusulkan WDCNN dengan kernel
konvolusi lebar (lebar 64, \emph{stride} 16) yang mampu memproses sinyal
vibrasi mentah secara langsung dan mencapai akurasi tinggi pada dataset
CWRU.
Kelebihan WDCNN adalah kemampuannya menangkap morfologi impuls transien
tanpa ekstraksi fitur manual.
Namun, interpretabilitas internal WDCNN---fitur mana dari sinyal yang
mendorong klasifikasi---belum dikaji secara sistematis dengan metode XAI
tingkat sinyal.
\textcolor{blue}{Sejumlah studi sebelumnya telah mencoba menerapkan XAI pada
diagnostik \emph{bearing}, namun beroperasi pada tingkat fitur agregat
atau resolusi spektrogram.}
Sebagai perbandingan, \textcolor{blue}{\citenameitb{Santos2024SHAPSVM}
menerapkan SHAP pada 9--15 fitur statistik untuk SVM,
sementara \citenameitb{Jeon2023CAMBearing} menggunakan
\emph{class activation mapping} (CAM) pada spektrogram STFT yang
beroperasi di resolusi frekuensi, bukan di resolusi sinyal mentah.}
Celah metodologis yang belum terisi adalah: belum ada studi yang menerapkan
SHAP DeepExplainer langsung pada sinyal mentah resolusi penuh (2.048 titik)
untuk WDCNN, memformalisasi peta atribusi per kelas (\emph{Fault Signature
Maps}), dan memvalidasinya dengan metrik kuantitatif yang sistematis.

\textbf{Prognostik RUL.}
Berbagai arsitektur \emph{deep learning}---LSTM \citetitb{Hochreiter1997LSTM},
Transformer \citetitb{Vaswani2017Attention}, Mamba \citetitb{Gu2023Mamba},
dan xLSTM \citetitb{Beck2024xLSTM} telah mencapai RMSE mendekati saturasi
pada dataset benchmark PHM2012 dan XJTU-SY.
Peningkatan akurasi arsitektural kini memberikan hasil yang semakin kecil
dibandingkan upaya rekayasa yang dibutuhkan.
Di sisi interpretabilitas, metode XAI yang banyak digunakan untuk RUL yaitu SHAP,
\emph{Integrated Gradients}, dan \emph{attention rollout}, beroperasi
di tingkat masukan. Metode tersebut mengidentifikasi langkah waktu mana yang
berpengaruh, namun tidak dapat menjawab apakah representasi laten model
bersesuaian dengan konsep fisik yang bermakna.
Mekanisme interpretabilitas (\emph{mechanistic interpretability}), yang
pertama kali dikembangkan untuk model bahasa besar dengan \emph{Sparse
Autoencoder} \textcolor{blue}{\citetitb{Bricken2023Monosemanticity}}, menawarkan pendekatan
berbeda: mengekstrak fitur monosemantik dari representasi laten yang
kemudian dapat dipetakan ke konsep fisik yang diketahui.
Ruang yang belum terisi adalah adaptasi pendekatan ini ke domain
prognostik \emph{bearing} belum pernah dilaporkan, termasuk prosedur
pemetaan ke frekuensi karakteristik degradasi dengan inferensi statistik
yang ketat.
Tanpa prosedur seperti ini, pertanyaan paling mendasar bagi insinyur
perawatan---\textcolor{blue}{\emph{``apakah model RUL ini benar-benar mengetahui tentang
BPFO?''}}---tidak dapat dijawab secara kuantitatif.
Ketersediaan jawaban yang terverifikasi atas pertanyaan ini merupakan
prasyarat untuk penggunaan model prediksi RUL berbasis \emph{deep
learning} di lingkungan industri yang kritis terhadap keselamatan.

\textbf{Konsistensi Berbagai Model.}
Sebagian besar studi diagnostik \emph{bearing} membandingkan satu arsitektur
dengan \emph{baseline} tunggal, sehingga tidak memberikan gambaran konsistensi
lintas golongan algoritma pada dataset dan protokol yang sama.
Selain itu, sebagian besar evaluasi menggunakan pembagian data acak
(\emph{random split}) yang dapat menghasilkan inflasi akurasi karena segmen
berurutan dari rekaman yang sama masuk ke set pelatihan dan pengujian
sekaligus, sebuah bentuk kebocoran data (\emph{data leakage}) yang sering
tidak dilaporkan dalam literatur.
Perbandingan sistematis antara golongan statistik-\emph{kernel} (SVM, LR),
\emph{tree-based} (DT, RF, XGBoost), dan \emph{deep learning} (WDCNN) dengan
protokol pemisahan temporal yang ketat dan inferensi statistik yang terstandarisasi
belum dilaporkan secara komprehensif.
Celah ini penting secara praktis karena industri membutuhkan panduan
yang dapat diandalkan, bukan sekadar perbandingan untuk
memilih model yang sesuai dengan sumber daya komputasi yang tersedia di setiap
tingkatan arsitektur sistem PdM.

\textcolor{blue}{Tabel~\ref{tab:bab1_posisi_penelitian} merangkum posisi
penelitian disertasi ini relatif terhadap literatur terdahulu pada sembilan
dimensi kritis yang teridentifikasi dari kajian \emph{state of the art}
di atas, dengan kontribusi penelitian disangga oleh dua artikel jurnal utama
(FSM dan SAE-BPFx) dan dua makalah konferensi pendukung.}

\begin{table}[!ht]
  \centering
  \caption{\textcolor{blue}{Posisi penelitian disertasi ini relatif terhadap
           literatur terdahulu pada sembilan dimensi kritis. Tanda kurung di
           kolom kontribusi menunjukkan sumber empiris internal:
           Conf~1 dan Conf~2 (konferensi pendukung), J~1 (FSM), J~2 (SAE-BPFx).}}
  \label{tab:bab1_posisi_penelitian}
  \scriptsize
  \setlength{\tabcolsep}{3pt}
  \renewcommand{\arraystretch}{1.2}
  \begin{tabularx}{\textwidth}{@{}>{\raggedright\arraybackslash}p{1.9cm}
                                   >{\raggedright\arraybackslash}p{3.0cm}
                                   >{\raggedright\arraybackslash}p{2.8cm}
                                   >{\raggedright\arraybackslash}X@{}}
    \toprule
    \textbf{Aspek} & \textbf{Penelitian Terdahulu} &
    \textbf{Gap / Keterbatasan} & \textbf{Kontribusi Disertasi Ini} \\
    \midrule

    Pendekatan model diagnostik &
    Model tunggal atau 2--3 algoritma: ANN+SVM
    \citetitb{Ayvaz2021IoTPdM}; SVM+k-NN
    \citetitb{Cui2022AnomalyML}; ANFIS
    \citetitb{Marani2020ANFIS} &
    Belum ada evaluasi komprehensif lintas golongan (ML tradisional vs.\
    \emph{deep learning}) pada dataset dan protokol evaluasi temporal
    yang sama &
    \emph{Benchmark} tiga golongan algoritma: statistik-kernel (SVM, LR),
    \emph{tree-based} (DT, RF, XGBoost), dan \emph{deep learning} (WDCNN)
    pada CWRU dengan \emph{temporal split} 54/13/33\,\%
    [Conf~1, Conf~2, J~1] \\[3pt]

    Pendekatan model prognostik &
    LSTM tunggal \citetitb{Hochreiter1997LSTM};
    Transformer \citetitb{Vaswani2017Attention};
    model tunggal per studi \citetitb{Lei2018PrognosticsReview} &
    Perbandingan lintas arsitektur modern (\emph{state-space},
    \emph{basis-block}, \emph{gated convolution}) pada dataset dan
    protokol identik belum dilaporkan &
    Tiga \emph{backbone} RUL: Mamba-xLSTM-Net, N-BEATS-xLSTM-RUL,
    SparseGate-TCN-RUL dilatih pada PHM2012, XJTU-SY, dan IMS dengan
    protokol identik (75 epoch, 3 \emph{seed}) [J~2] \\[3pt]

    XAI diagnostik (level sinyal) &
    XAI pada fitur agregat: SHAP+SVM pada 9--15 fitur statistik
    \citetitb{Santos2024SHAPSVM}; CAM pada spektrogram STFT
    \citetitb{Jeon2023CAMBearing} &
    Belum ada XAI pada sinyal vibrasi mentah resolusi penuh
    (2.048 titik) untuk CNN \emph{bearing} dengan validasi
    multi-metrik &
    \emph{Fault Signature Maps} (FSM): SHAP DeepExplainer pada sinyal mentah
    2.048 titik, tiga varian (Signed, Absolute, Variance), tiga metrik
    kuantitatif --- \emph{discriminability} $D = 0{,}216$,
    \emph{severity monotonicity}, \emph{split-half stability} $r = 0{,}940$
    --- plus ablasi \emph{BatchNorm} (+133\,\% diskriminabilitas dengan
    \emph{trade-off} akurasi $-3{,}60$~pp) [J~1] \\[3pt]

    XAI prognostik (level laten) &
    XAI di tingkat masukan: SHAP, \emph{Integrated Gradients},
    \emph{attention rollout} \citetitb{Serradilla2022DLPdM} &
    Tidak dapat menjawab apakah representasi laten model RUL
    bersesuaian dengan fisika degradasi \emph{bearing}
    (BPFO, BPFI, BSF, FTF) &
    SAE-BPFx: \emph{Top-k Sparse Autoencoder} dipetakan ke frekuensi
    karakteristik BPFx melalui korelasi Pearson dengan spektrum
    \emph{Hilbert envelope};     \emph{hit-rate} BPFI 2{,}3\,\% pada PHM2012,
    BSF 1{,}6\,\% pada XJTU-SY ($p < 0{,}001$) [J~2] \\[3pt]

    Validasi statistik &
    Akurasi dan F1 saja tanpa inferensi statistik
    \citetitb{Serradilla2022DLPdM}; sebagian besar tanpa CI atau
    uji hipotesis &
    Klaim kinerja tidak didukung oleh interval kepercayaan,
    uji hipotesis, atau kontrol negatif &
    \emph{Bootstrap} CI 95\,\% ($B = 1.000$), \emph{permutation test}
    dua sisi ($B = 1.000$), dua \emph{negative controls} (Xavier-init,
    \emph{Gaussian noise}); tiga metrik FSM untuk diagnostik [J~1, J~2] \\[3pt]

    Konsistensi lintas arsitektur prognostik &
    Studi tunggal per arsitektur tanpa uji universalitas
    \citetitb{Lei2018PrognosticsReview} &
    Tidak dapat dipastikan apakah pemetaan BPFx merupakan properti
    universal representasi terlatih atau artefak arsitektur tertentu &
    Urutan dominansi BPFx (BPFI~$\succ$~BPFO~$\succ$~BSF~$\succ$~FTF)
    konsisten lintas tiga \emph{backbone} (Mamba/N-BEATS/SparseGate)
    meski dengan variasi magnituda 4--8 kali lipat [J~2] \\[3pt]

    Dataset dan validasi industri &
    Dataset publik saja (CWRU, IMS, PHM) tanpa validasi industri
    \citetitb{Serradilla2022DLPdM,Khorram2020CNNLSTM} &
    Celah antara performa akademik dan penerapan industri masih besar;
    validasi pada data produksi nyata sangat terbatas &
    Empat dataset publik (CWRU, PHM2012, XJTU-SY, IMS) plus validasi
    lintas domain RUL pada data vibrasi real mesin \emph{grinding}
    OR1/OR2 PT~SKF Indonesia \citetitb{SKF2017TechReport} \\[3pt]

    Arsitektur sistem PdM &
    \emph{Cloud-centric} \citetitb{Cheng2019EdgeCloud};
    \emph{edge} tanpa lapisan XAI \citetitb{Hafeez2021EdgePdM} &
    Belum ada arsitektur \emph{multi-tier} yang mengintegrasikan diagnostik
    dan prognostik beserta XAI per lapis secara terstruktur &
    Arsitektur tiga-\emph{tier}: Edge IoT (fitur statistik) $\to$
    Edge Server (WDCNN+FSM sebagai lapisan audit) $\to$ Cloud
    (\emph{backbone} RUL dengan SAE-BPFx) sebagai cetak biru adopsi
    industri \\[3pt]

    Integrasi data kualitas produk &
    PdM dan kontrol kualitas dianggap dua sistem terpisah; belum ada
    umpan balik tertutup antara hasil pemeriksaan akhir (\emph{vibration
    checking}, \emph{radial clearance checking}) dan keputusan
    perawatan &
    Data kualitas produk tidak dimanfaatkan sebagai sinyal dini
    kerusakan; masalah kualitas terdeteksi setelah produksi tidak
    berkontribusi pada perawatan prediktif &
    Peta jalan integrasi data kualitas produk (\emph{vibration checking},
    \emph{radial clearance checking}) PT~SKF Indonesia sebagai aliran
    data tambahan untuk model multi-modal pada penelitian lanjutan;
    arsitektur \emph{multi-tier} disiapkan untuk mengakomodasi aliran
    ini di lapisan Cloud \\

    \bottomrule
  \end{tabularx}

  \smallskip
  {\scriptsize Sumber empiris internal:
   Conf~1 = \citenameitb{TotoSuharto2024Conf1SVM};
   Conf~2 = \citenameitb{TotoSuharto2024Conf2Tree};
   J~1 = \citenameitb{TotoSuharto2025Journal1FSM};
   J~2 = \citenameitb{TotoSuharto2025Journal2RUL}.}
\end{table}

\textcolor{blue}{Pemaparan detail \emph{state of the art} per studi, termasuk
metode, dataset, hasil utama, dan keterbatasan masing-masing, disajikan pada
Tabel~\ref{tab:bab2_sota_pdm} di Bab~\ref{bab:tinjauan-pustaka}.}


%% ================================================================
\section{Tujuan Penelitian}
\label{sec:bab1_tujuan}
%% ================================================================

Tujuan penelitian pada disertasi ini adalah menghasilkan kerangka perawatan prediktif
\emph{bearing} yang tidak hanya akurat, tetapi juga dapat diinterpretasikan dan dapat dipercaya oleh insinyur perawatan di lingkungan industri
manufaktur\textcolor{blue}{, sekaligus menyediakan fondasi konseptual untuk
integrasi data sensor mesin dengan informasi kualitas produk secara waktu nyata}.
Untuk mencapai tujuan tersebut, diperlukan tujuan-tujuan yang lebih rinci
sebagai berikut:

\begin{enumerate}

  \item \textcolor{blue}{Memvalidasi dan membandingkan kinerja beberapa golongan
        algoritma diagnostik \emph{bearing} (statistik-\emph{kernel},
        \emph{tree-based}, dan \emph{deep learning}) serta mengembangkan model prognostik RUL
        sehingga tersedia model diagnostik dan prognostik yang akurat
        sebagai fondasi sistem perawatan prediktif \emph{ball bearing}.}

  \item \textcolor{blue}{Mengembangkan kerangka XAI untuk model
        perawatan prediktif \emph{ball bearing}: (a)~pada level diagnostik,
        menerapkan SHAP DeepExplainer pada sinyal vibrasi mentah resolusi
        penuh (2.048 titik) untuk memformalisasi \emph{Fault Signature Maps}
        (FSM) yang divalidasi dengan metrik \emph{discriminability index},
        \emph{severity monotonicity}, dan \emph{split-half stability}; dan
        (b)~pada level prognostik, mengadaptasi \emph{Top-k Sparse
        Autoencoder} (SAE) untuk memetakan representasi laten model RUL ke
        frekuensi karakteristik BPFO, BPFI, BSF, dan FTF melalui
        \emph{bootstrap} CI dan \emph{permutation test}.}

  \item \textcolor{blue}{Menerapkan dan memvalidasi model prediksi RUL
        \emph{bearing} yang dikembangkan pada dataset benchmark terhadap data
        sensor vibrasi dari proses produksi \emph{ball bearing} kondisi nyata,
        serta merancang arsitektur sistem perawatan
        prediktif multi-tier yang mengintegrasikan komponen diagnostik (FSM) dan prognostik
        (SAE-BPFx) untuk adopsi industri dan pengurangan waktu henti
        yang tidak terencana pada produksi \emph{ball bearing}.}

\end{enumerate}


%% ================================================================
\section{Manfaat Penelitian}
\label{sec:bab1_manfaat}
%% ================================================================

\textcolor{blue}{Penelitian ini diharapkan memberikan manfaat sebagai berikut:}

\begin{enumerate}

      \item \textcolor{blue}{Memberikan contoh kasus implementasi analisis
      big data dan AI dalam sistem produksi dalam meningkatkan produktivitas dan daya saing.}

      \item \textcolor{blue}{Secara ilmiah, penelitian ini berkontribusi pada
        pengembangan metode XAI yang beroperasi di level sinyal mentah
        (diagnostik) dan level representasi laten (prognostik) untuk perawatan
        prediktif \emph{bearing}, yang selama ini masih jarang dilakukan.
        Hasil penelitian berupa kerangka FSM dan SAE-BPFx memperkaya
        literatur interpretabilitas model \emph{deep learning} pada domain
        mesin rotasi.}

  \item \textcolor{blue}{Secara praktis, validasi model pada data vibrasi
        industri real menunjukkan kelayakan penerapan
        sistem perawatan prediktif berbasis \emph{deep learning} untuk
        mengurangi waktu henti yang tidak terencana di lingkungan manufaktur.
        Arsitektur multi-tier yang dirancang memberikan cetak biru yang dapat
        diadopsi oleh industri untuk implementasi bertahap.}

%   \item \textcolor{blue}{Hasil penelitian yang berupa \emph{pipeline} sumber
%         terbuka, bobot model, dan kerangka konseptual multi-tier dapat menjadi
%         contoh dan acuan bagi implementasi perawatan prediktif berbasis
%         kecerdasan artifisial pada sistem produksi lain yang menggunakan
%         proses permesinan, tidak terbatas pada produksi \emph{ball bearing}.}
   \item \textcolor{blue}{Hasil dari penelitian ini yang berupa sistem 
   dengan menggunakan perangkat lunak dan perangkat keras dapat 
   menjadi contoh bagi implementasi analisis big data dan AI dalam memperbaiki 
   sistem perawatan prediktif yang sudah ada.}

\end{enumerate}


%% ================================================================
\section{Kontribusi dan Kebaruan}
\label{sec:bab1_kontribusi}
%% ================================================================

Penelitian disertasi yang telah dilakukan memberikan kontribusi pada
beberapa aspek berikut:

\begin{enumerate}

  \item Penyediaan kerangka interpretabilitas dari masukan sinyal mentah (\emph{signal-level
        XAI}) yang memungkinkan verifikasi keputusan model
        \emph{deep learning} untuk diagnostik \emph{bearing}.
        Berbeda dari pendekatan XAI berbasis fitur agregat yang selama ini
        mendominasi literatur, kerangka FSM beroperasi langsung pada 2.048
        titik sinyal mentah sehingga pola atribusi dapat dipetakan ke
        morfologi impuls transien yang dikenal dari teori kerusakan
        \emph{bearing} \citetitb{TotoSuharto2025Journal1FSM}.

  \item Pengembangan prosedur pemetaan kuantitatif yang menghubungkan
        representasi laten model prediksi RUL \emph{bearing} dengan
        frekuensi karakteristik fisik degradasi, dilengkapi
        \emph{bootstrap} CI 95\,\%, \emph{permutation test}, dan
        \emph{negative controls} untuk memastikan korespondensi adalah
        properti representasi terlatih, bukan artefak arsitektur
        \citetitb{TotoSuharto2025Journal2RUL}.

  \item Penyediaan \emph{benchmark} multi-model yang sistematis dengan
        inferensi statistik ketat, membandingkan tiga \textcolor{blue}{golongan} algoritma
        (statistik-\emph{kernel}, \emph{tree-based}, \emph{deep learning})
        dengan evaluasi temporal yang sama.
        Hasilnya memberikan panduan pemilihan model yang dapat diandalkan
        untuk adopsi industri dengan berbagai tingkatan komputasi
        \citetitb{TotoSuharto2024Conf1SVM,TotoSuharto2024Conf2Tree}.

  \item \textcolor{blue}{Validasi lintas domain model prediksi RUL pada data
        sensor vibrasi dari proses produksi \emph{ball bearing} di
        PT~SKF~Indonesia, yang menunjukkan kelayakan penerapan pipeline
        prognostik yang dikembangkan pada dataset benchmark ke lingkungan
        industri manufaktur nyata, disertai perancangan arsitektur multi-tier
        yang menjadi cetak biru untuk adopsi industri dan peta jalan bagi
        integrasi data kualitas produk pada penelitian lanjutan.}

  \item Penyediaan \emph{pipeline} sumber terbuka---\emph{notebook}
        reproduksi, bobot model, dan skrip evaluasi statistik---yang
        dapat diadopsi oleh peneliti lain untuk domain mesin rotasi
        berbeda, mempercepat kemajuan penelitian di bidang PdM.

\end{enumerate}

Secara lebih rinci, beberapa hasil penelitian yang telah dilaksanakan
beserta kebaruannya dijelaskan sebagai berikut.

\textbf{1. \emph{Fault Signature Maps} (FSM): XAI tingkat sinyal mentah
untuk WDCNN.}
Ini merupakan penerapan SHAP DeepExplainer secara langsung pada
sinyal vibrasi mentah resolusi penuh dengan 2.048 titik per segmen per kelas untuk
arsitektur WDCNN pada dataset CWRU.
FSM diformalisasikan dalam tiga varian matematis (Signed, Absolute, Variance)
dan divalidasi melalui tiga metrik kuantitatif: \emph{discriminability
index} $D = 0{,}216$, \emph{severity monotonicity} Ball 17{,}6\,\% /
IR 13{,}8\,\% / OR 8{,}6\,\%, dan \emph{split-half cosine similarity}
\textcolor{blue}{$r = 0{,}940$ pada analisis utama (300 \emph{background} /
500 sampel uji)}.
Metode ablasi menjelaskan \emph{trade-off} yang belum pernah dilaporkan
sebelumnya: \textcolor{blue}{pada subset SHAP tereduksi (30 \emph{background}
/ 50 sampel uji), menghapus \emph{BatchNorm} meningkatkan
\emph{discriminability} FSM sebesar 133\,\% (dari $0{,}315$ menjadi $0{,}735$)}
dengan \emph{trade-off} akurasi $-3{,}60$~pp,
memberikan panduan desain arsitektur antara akurasi dan interpretabilitas
\citetitb{TotoSuharto2025Journal1FSM}.

\textbf{2. SAEBearing: Interpretabilitas mekanistik untuk prognostik RUL.}
Ini merupakan adaptasi pertama \emph{Top-k Sparse Autoencoder}
(\emph{mechanistic interpretability}), yang sebelumnya dikembangkan
untuk model bahasa besar, ke domain prognostik RUL \emph{bearing}.
SAE dilatih secara \emph{post-hoc} pada representasi laten tiga arsitektur
\emph{backbone} RUL---Mamba-xLSTM, N-BEATS-xLSTM, SparseGate-TCN---dan
fitur-fitur latennya dipetakan ke frekuensi karakteristik \emph{bearing}
melalui korelasi Pearson dengan spektrum \emph{Hilbert envelope}.
Prosedur falsifikasi menggunakan \emph{bootstrap} CI 95\,\%,
\emph{permutation test}, dan dua \emph{negative controls} (inisialisasi
Xavier dan \emph{Gaussian noise}) untuk memastikan korespondensi adalah
properti representasi terlatih, bukan artefak arsitektur.
Hasil utama: \emph{hit-rate} BPFI 2{,}3\,\% pada PHM2012
($p < 0{,}001$); BSF 1{,}6\,\% pada XJTU-SY ($p < 0{,}001$);
BPFI 1{,}76\,\% dan BSF 0{,}49\,\% pada IMS; korelasi tertinggi
$r_{\max} = 0{,}507$.
Urutan dominansi BPFx konsisten lintas tiga arsitektur \emph{backbone}
meski dengan variasi magnituda 4--8 kali lipat
\citetitb{TotoSuharto2025Journal2RUL}.

\textbf{3. \textcolor{blue}{Validasi lintas domain dan kerangka PdM multi-tier.}}
\textcolor{blue}{Model prediksi RUL yang dikembangkan pada dataset benchmark
(PHM2012, XJTU-SY, IMS) divalidasi pada data sensor vibrasi dari mesin
\emph{grinding} OR1/OR2 PT~SKF~Indonesia.
Validasi lintas domain ini menunjukkan kelayakan penerapan pipeline prognostik
yang dikembangkan pada kondisi laboratorium ke lingkungan produksi
\emph{ball bearing} yang sesungguhnya, dengan kerusakan yang berkembang
secara alami.}
Sintesis komponen FSM (\emph{input-attribution}, \textcolor{blue}{diagnostik}) dan SAE-BPFx
(\emph{latent-concept}, \textcolor{blue}{prognostik}) ke dalam cetak biru sistem perawatan
prediktif tiga tingkatan: (i)~Edge IoT untuk triage anomali ringan dengan
model ringan berbasis fitur statistik, (ii)~Edge Server untuk klasifikasi
kerusakan akurat menggunakan WDCNN dengan FSM sebagai lapisan audit, dan
(iii)~Cloud untuk estimasi RUL jangka panjang menggunakan arsitektur
\emph{backbone} dengan interpretabilitas SAE-BPFx.
\textcolor{blue}{Kerangka ini diperkuat oleh hasil validasi pada data
industri real PT~SKF~Indonesia, dan menyertakan peta jalan integrasi data
kualitas produk (\emph{vibration checking} dan \emph{radial clearance
checking}) sebagai aliran data tambahan untuk penelitian lanjutan.}
Arsitektur multi-tier bersifat konseptual---sebagai cetak biru---namun
validasi RUL pada data industri real memberikan bukti awal kelayakan
\emph{deployment}.


%% ================================================================
\section{Metodologi Penelitian}
\label{sec:bab1_metodologi}
%% ================================================================

Penelitian disertasi ini telah melalui tahap-tahap penelitian sesuai
dengan standar pelaksanaan penelitian doktoral di bidang \textcolor{blue}{teknik
dan manajemen industri}.
Metodologi penelitian dirancang mengikuti enam tahap makro berikut,
kemudian dijabarkan secara operasional pada rincian pelaksanaan di akhir
subbab ini.

\textcolor{blue}{
\subsection{Tahap Perencanaan}
\label{subsec:bab1_tahap_perencanaan}

Tahap perencanaan dimulai dengan identifikasi permasalahan melalui studi
literatur, perumusan \emph{state of the art}, rumusan masalah dan tujuan
penelitian, penetapan ruang lingkup (batasan dan asumsi), serta
pemahaman objek kajian.
Kajian literatur difokuskan pada penerapan \emph{big data}, \emph{machine
learning}, \emph{deep learning}, dan \emph{explainable AI} dalam
perawatan prediktif \emph{bearing}.
Celah teori (\emph{gap}) yang teridentifikasi---interpretabilitas sinyal
mentah pada diagnostik dan interpretabilitas mekanistik laten pada
prognostik---menjadi landasan perumusan tiga butir rumusan masalah pada
Subbab~\ref{sec:bab1_rumusan}.
Posisi penelitian dipetakan terhadap literatur SoTA yang telah diuraikan
pada Bab~\ref{bab:tinjauan-pustaka}, sementara skema pendekatan keseluruhan
diilustrasikan pada Gambar~\ref{fig:bab1_skema}.
Objek empiris utama meliputi dataset \emph{benchmark} publik (CWRU,
PHM2012, XJTU-SY, IMS) dan validasi terbatas pada data vibrasi industri
PT~SKF~Indonesia.

\subsection{Rancangan dan Persiapan Penelitian}
\label{subsec:bab1_rancangan_persiapan}

Pada tahap ini dilakukan perancangan dan persiapan penelitian sebelum
pelatihan model.
Langkah pertama adalah pengumpulan dan kurasi data penelitian dari
dataset publik dan sumber industri, disertai dokumentasi protokol
pemisahan \emph{bearing}-wise agar tidak terjadi kebocoran informasi
antar set pelatihan, validasi, dan uji.
Langkah kedua adalah menyusun kerangka praproses dan rekayasa fitur yang
konsisten untuk jalur diagnostik dan prognostik.
Kerangka tersebut mengikuti alur berikut.
}

\textcolor{blue}{
\begin{enumerate}
  \item \textbf{Ekstraksi fitur (\emph{feature generation}).}
        Data mentah sinyal getaran diubah menjadi representasi yang siap
        dimodelkan: segmen jendela geser untuk diagnostik CWRU, vektor
        HI~36-D untuk model berbasis kernel/pohon, serta tensor jendela
        HI untuk model \emph{deep learning} prognostik.
        Detail teknis diuraikan pada Bab~\ref{bab:metodologi-umum}
        Subbab~\ref{sec:bab3_praproses}.
  \item \textbf{Pembersihan dan normalisasi fitur (\emph{feature cleaning
        and normalization}).}
        Rekaman dengan artefak akuisisi, nilai hilang, duplikat, dan
        \emph{outlier} ditangani sebelum pemodelan; normalisasi
        Z-score/\emph{StandardScaler} dikalibrasi hanya pada data
        pelatihan.
  \item \textbf{Penskalaan fitur (\emph{feature scaling}).}
        Fitur dari ranah waktu dan frekuensi diseragamkan skalanya agar
        tidak ada fitur yang dominan semata karena rentang numeriknya
        lebih besar.
  \item \textbf{Seleksi fitur (\emph{feature selection}).}
        Subset fitur relevan dipilih untuk model berbasis kernel dan
        pohon keputusan; model \emph{deep learning} menggunakan
        representasi tersegmentasi langsung.
  \item \textbf{Transformasi data (\emph{data transformation}).}
        Pelabelan RUL (\emph{piecewise} atau linear), segmentasi temporal,
        dan penyiapan \emph{batch} pelatihan disesuaikan per dataset
        (PHM2012, XJTU-SY, IMS, CWRU).
\end{enumerate}
}

\textcolor{blue}{
Dua kelompok tugas pemodelan dirancang paralel.
Untuk \textbf{regresi} (prognostik RUL), metrik evaluasi utama adalah
RMSE, MAE, $R^2$, dan PHM~Score (PHM2012).
Untuk \textbf{klasifikasi} (diagnostik CWRU), metrik utama adalah akurasi,
\emph{precision}, \emph{recall}, dan F1.
Perbaikan model (\emph{model improvement}) meliputi \emph{hyperparameter
tuning}, validasi silang (\emph{k-fold cross-validation} pada jalur
klasifikasi), pelatihan multi-\emph{seed} (42, 43, 44) pada jalur
prognostik, dan protokol \emph{fixed-budget training} (75 epoch,
\texttt{bf16-mixed}).
}

\textcolor{blue}{
\subsection{Pengembangan Model}
\label{subsec:bab1_pengembangan_model}

Pengembangan model dilakukan berdasarkan hasil \emph{benchmark} awal dan
karakteristik masing-masing dataset.
Pada jalur diagnostik, tiga golongan algoritma (statistik-\emph{kernel},
\emph{tree-based}, \emph{deep learning}) dibandingkan sebelum WDCNN dan
kerangka FSM dikembangkan.
Pada jalur prognostik, dua belas arsitektur disaring melalui praskrining
30~epoch; tiga \emph{backbone} Tier-S (Mamba-xLSTM-Net, N-BEATS-xLSTM-RUL,
SparseGate-TCN-RUL) dilatih 75~epoch dengan tiga \emph{seed}.
Adaptasi \emph{Top-k Sparse Autoencoder} (\emph{SAEBearing}) dilakukan
\emph{post-hoc} pada \emph{hidden states} model terbaik untuk pemetaan
BPFx.
}

\textcolor{blue}{
\subsection{Komparasi Model dengan Dunia Nyata}
\label{subsec:bab1_komparasi_dunia_nyata}

Komparasi model dengan dunia nyata merupakan langkah validasi untuk
menguji sejauh mana model yang dikembangkan pada dataset \emph{benchmark}
dapat digeneralisasi.
Validasi lintas domain dilakukan dengan menerapkan model prognostik terbaik
pada data vibrasi mesin \emph{grinding} OR1/OR2 PT~SKF~Indonesia.
Komparasi sebelum--sesudah pengembangan model dievaluasi melalui metrik
RUL dan analisis residual per \emph{bearing}.
}

\textcolor{blue}{
\subsection{Aksi untuk Perbaikan Situasi}
\label{subsec:bab1_aksi_perbaikan}

Tahap aksi perbaikan situasi mensintesiskan temuan diagnostik dan
prognostik ke dalam kerangka perawatan prediktif multi-\emph{tier}
(Edge~IoT --- Edge~Server --- Cloud/GPU) yang dapat diadopsi pada
lini produksi \emph{ball bearing}.
Kerangka ini menghubungkan deteksi kerusakan berbasis FSM, estimasi RUL
berbasis \emph{backbone} Tier-S, dan interpretabilitas mekanistik
SAE-BPFx ke dalam alur keputusan perawatan yang dapat diaudit insinyur
pabrik.
}

\textcolor{blue}{
\subsection{Tahap Kesimpulan dan Saran Penelitian}
\label{subsec:bab1_tahap_kesimpulan}

Tahap penutup penelitian merumuskan jawaban terhadap tujuan penelitian
berdasarkan bukti empiris Bab~IV dan Bab~V, mengidentifikasi
keterbatasan metodologis, serta memberikan rekomendasi penelitian
lanjutan---termasuk integrasi data kualitas produk \emph{real-time} dan
perluasan validasi industri.
Uraian lengkap disajikan pada Bab~\ref{bab:kesimpulan}.
}

\textcolor{blue}{Rincian operasional pelaksanaan metodologi dijabarkan
pada poin-poin berikut.}

\begin{enumerate}

  \item \textbf{Studi Literatur.}
        Kajian literatur dilakukan secara komprehensif terhadap lebih dari
        seratus makalah di bidang perawatan prediktif, diagnostik
        \emph{bearing}, prognostik RUL, dan \emph{explainable AI} untuk
        mesin rotasi.
        Sebagian besar makalah bersumber dari jurnal terkemuka seperti
        IEEE Transactions, \emph{Mechanical Systems and Signal Processing},
        dan \emph{Reliability Engineering \& System Safety}.
        Hasil kajian literatur memetakan \emph{state of the art},
        mengidentifikasi celah penelitian yang belum terisi, dan
        memotivasi \textcolor{blue}{kedua komponen} penelitian utama (FSM dan SAE-BPFx).

  \item \textbf{Persiapan Data Penelitian.}
        Penelitian menggunakan empat dataset \emph{bearing} yang bersifat
        publik.
        Untuk diagnostik: dataset CWRU (\emph{Case Western Reserve
        University}) dengan sinyal \emph{drive-end} pada frekuensi sampling
        48~kHz, sepuluh kelas kerusakan (Normal + IR/OR/Ball pada tiga
        ukuran kerusakan), segmen 2.048 titik, dan pembagian temporal
        54/13/33\,\%.
        Untuk prognostik: dataset PHM2012 (FEMTO-PRONOSTIA, 17
        \emph{bearing}), XJTU-SY (LDK~UER204, 10~\emph{bearing}), dan
        IMS (Rexnord, 4~\emph{bearing}).
        Praproses meliputi normalisasi dengan StandardScaler yang
        dikalibrasi hanya pada data pelatihan untuk mencegah kebocoran
        informasi.

  \item \textbf{Penelitian Awal: Benchmark Multi-Model.}
        Penelitian awal dilakukan dengan melatih dan mengevaluasi tiga
        golongan algoritma diagnostik \emph{bearing} secara paralel pada
        dataset CWRU: golongan statistik-\emph{kernel} (SVM-RBF, LR),
        \emph{tree-based} (DT, RF, XGBoost), dan \emph{deep learning}
        (WDCNN).
        Evaluasi dilakukan dengan protokol pembagian temporal yang ketat
        (bukan pembagian acak) untuk menghindari inflasi akurasi akibat
        kebocoran data antar-segmen.
        Hasil \emph{benchmark} ini memberikan dasar konsistensi lintas
        golongan algoritma dan motivasi untuk lapisan XAI per golongan.

  \item \textbf{Pengembangan Kerangka FSM.}
        Kerangka \emph{Fault Signature Maps} (FSM) dikembangkan melalui
        penerapan SHAP DeepExplainer pada model WDCNN yang telah dilatih,
        dengan 300 sampel \emph{background} dan 500 sampel uji yang
        diambil secara stratifikasi.
        Hasilnya adalah 10 array atribusi berukuran $500 \times 2.048$,
        menghasilkan lebih dari 10 juta nilai SHAP individual.
        FSM diformalisasikan dalam tiga varian dan divalidasi dengan tiga
        metrik kuantitatif, dilengkapi studi ablasi tiga varian arsitektur
        untuk mengungkap \emph{trade-off} antara akurasi dan
        interpretabilitas.

  \item \textbf{Pengembangan SAEBearing.}
        Adaptasi \emph{Top-k Sparse Autoencoder} ke domain prognostik RUL
        dilakukan melalui protokol tiga langkah: (i)~pelatihan tiga
        arsitektur \emph{backbone} RUL pada masing-masing dataset
        (\emph{fixed-budget training}, 75 epoch, 3 \emph{seed}),
        (ii)~pelatihan SAE \emph{post-hoc} pada kumpulan \emph{hidden
        states} dari model yang telah beku (\emph{frozen}), dan
        (iii)~korelasi fitur laten SAE dengan spektrum \emph{Hilbert
        envelope} pada frekuensi karakteristik yang dihitung dari geometri
        \emph{bearing} masing-masing dataset.

  \item \textbf{Pengujian dan Validasi Statistik.}
        Seluruh klaim kontribusi divalidasi dengan inferensi statistik yang
        ketat: \emph{bootstrap} CI 95\,\% (\textcolor{blue}{$B = 1.000$ iterasi}), \emph{permutation
        test} (\textcolor{blue}{$B = 1.000$}), dan dua \emph{negative controls} (inisialisasi
        Xavier dan \emph{Gaussian noise}) untuk SAEBearing.
        Uji konsistensi lintas arsitektur dilakukan dengan membandingkan
        urutan dominansi BPFx pada ketiga \emph{backbone} secara bersamaan.

  \item \textbf{Pemartifisial Laporan dan Publikasi.}
        Hasil penelitian dilaporkan dalam buku disertasi ini dan telah
        dipublikasikan atau sedang dalam proses \emph{review} pada jurnal
        dan konferensi internasional.
        Seluruh kode sumber, bobot model, dan \emph{notebook} reproduksi
        tersedia sebagai \emph{pipeline} sumber terbuka untuk mendukung
        reproduksi dan pengembangan lanjutan oleh peneliti lain.

\end{enumerate}

Keseluruhan tahapan metodologi di atas dirancang secara iteratif: temuan
dari satu tahap---misalnya identifikasi celah interpretabilitas pada hasil
\emph{benchmark}---menjadi masukan yang memperketat rancangan tahap
berikutnya.
Pendekatan ini memastikan bahwa setiap kontribusi dalam disertasi ini
terhubung secara logis dari satu ke berikutnya dan dapat direproduksi
secara penuh.


%% ================================================================
\section{Batasan dan Asumsi Penelitian}
\label{sec:bab1_batasan}
%% ================================================================

Penelitian ini dibatasi pada hal-hal berikut: (1)~sistem produksi yang
menggunakan proses permesinan \emph{grinding} pada produksi
\emph{rolling element bearing}; (2)~sinyal masukan adalah data vibrasi
mentah dari akselerometer, dengan fitur frekuensi hanya untuk validasi
fisika; (3)~penggunaan dataset diagnostik CWRU dan dataset prognostik PHM2012,
XJTU-SY, serta IMS; \textcolor{blue}{(4)~validasi lintas domain pada data
industri menggunakan data vibrasi dari mesin \emph{grinding} OR1/OR2
PT~SKF~Indonesia, dengan fokus pada komponen prognostik (prediksi RUL);}
(5)~integrasi data kualitas produk
(\emph{vibration checking} dan \emph{radial clearance checking})
secara waktu nyata merupakan peta jalan penelitian lanjutan, bukan
bagian dari implementasi empiris disertasi ini; dan (6)~seluruh
perangkat lunak menggunakan \emph{library} sumber terbuka yang
telah ada.

\textcolor{blue}{Adapun asumsi yang digunakan dalam penelitian ini adalah
sebagai berikut:}
(1)~sensor vibrasi telah terpasang dan terkalibrasi sesuai spesifikasi pada
mesin produksi \emph{ball bearing} (OR1 dan OR2);
(2)~mesin produksi beroperasi dalam jangkauan operasi normal sesuai dengan
spesifikasi teknis yang ditetapkan;
(3)~infrastruktur komputasi yang memadai (GPU) tersedia untuk pelatihan
model \emph{machine learning} dan \emph{deep learning};
\textcolor{blue}{(4)~data sensor vibrasi dari mesin \emph{grinding}
PT~SKF~Indonesia yang digunakan untuk validasi lintas domain telah diakuisisi
melalui \emph{gateway} SKF~IMx-8 sesuai protokol akuisisi standar}; dan
(5)~proses permesinan \emph{grinding} pada \emph{ball bearing} memiliki
karakteristik degradasi yang dapat dimodelkan menggunakan pendekatan
\emph{machine learning} berbasis data.


%% ================================================================
\section{Sistematika Penulisan}
\label{sec:bab1_sistematika}
%% ================================================================

Disertasi ini disusun ke dalam enam bab dengan alur yang koheren dari
konteks hingga kesimpulan.

\textbf{Bab~\ref{bab:pendahuluan}} (Pendahuluan) memuat latar belakang
urgensi perawatan prediktif dalam konteks \emph{Making Indonesia 4.0},
\textcolor{blue}{rumusan masalah tiga butir, tujuan penelitian, manfaat
penelitian}, kontribusi dan kebaruan beserta uraian teknis, metodologi
penelitian, batasan dan asumsi penelitian, serta sistematika penulisan.

\textbf{Bab~\ref{bab:tinjauan-pustaka}} (Tinjauan Pustaka) memetakan
\emph{state of the art} di bidang perawatan prediktif dan Industri~4.0,
diagnostik \emph{bearing} berbasis \emph{deep learning}, metode XAI
tingkat atribusi masukan, prognostik RUL, \emph{mechanistic
interpretability} dan \emph{Sparse Autoencoder}, fisika vibrasi
\emph{bearing}, serta peta celah literatur yang memotivasi penelitian ini.

\textbf{Bab~\textcolor{blue}{\ref{bab:metodologi-umum}}} (Metodologi) membahas metodologi penelitian yang digunakan pada
 Bab~\textcolor{blue}{\ref{bab:diagnostik}} dan Bab~\textcolor{blue}{\ref{bab:prognostik}}. Pembahasannya terkait kerangka konseptual terintegrasi, spesifikasi
empat dataset, praproses dan ekstraksi fitur, metrik evaluasi bersama,
serta infrastruktur dan protokol reproduksi.

\textbf{Bab~\textcolor{blue}{\ref{bab:diagnostik}}} (\emph{Bearing Fault Detection})
membahas penelitian terkait \textcolor{blue}{diagnostik} secara \emph{end-to-end} dan \emph{benchmark}
model algoritma dengan SHAP, formalisasi FSM dalam tiga
varian matematis, hasil klasifikasi pada dataset CWRU, studi ablasi
\emph{BatchNorm}, dan interpretasi fisika sinyal.
Bab ini disusun dengan urutan metode $\rightarrow$ hasil agar pembaca
dapat mengikuti alur dari rancangan eksperimen hingga temuan.

\textbf{Bab~\textcolor{blue}{\ref{bab:prognostik}}} (\emph{Remaining Useful Life})
menyajikan penelitian terkait \textcolor{blue}{ prognostik} secara \emph{end-to-end} dengan tiga arsitektur
\emph{backbone} RUL, \emph{Top-k Sparse Autoencoder}, pemetaan BPFx,
\emph{hit-rate} dengan \emph{bootstrap} CI dan \emph{permutation test},
\emph{negative controls}, dan perbandingan lintas arsitektur.
Bab ini disusun paralel dengan Bab~\textcolor{blue}{\ref{bab:diagnostik}} sehingga pembaca dapat
membandingkan pendekatan diagnostik dan prognostik secara simetris.

\textbf{Bab~\ref{bab:kesimpulan}} (Kesimpulan dan Rekomendasi) menutup
disertasi dengan pembahasan terintegrasi, rekapitulasi metodologi dan hasil, keterbatasan penelitian, dan
rekomendasi lanjutan, termasuk prioritas integrasi data \emph{real-world}
PT~SKF sebagai peta jalan konkret bagi peneliti selanjutnya.
